% =========================================================================
% Chapter 6: Conclusion and Future Prospects
% =========================================================================

\chapter{Conclusion and Future Prospects}\label{ch:conclusion}

The three objectives were met within the limits of a zone-of-inhibition screen.

Plant material of \taxa{Dipsacus inermis} was collected from the Kashmir University Botanical Garden at 34.12981$^{\circ}$\,N, 74.83396$^{\circ}$\,E (WGS84), authenticated at the Centre for Biodiversity and Taxonomy, dried, powdered, and extracted with distilled water in a Soxhlet apparatus. The concentrate was reconstituted at \SI{50}{\milli\gram\per\milli\litre} (C1) and \SI{100}{\milli\gram\per\milli\litre} (C2).

Qualitative tubes were positive for alkaloids, flavonoids, phenols, tannins, saponins, terpenoids, phytosterols, carbohydrates, and quinones.

Both extract concentrations produced zones against \taxa{Escherichia coli}, \taxa{Pseudomonas aeruginosa}, and \taxa{Klebsiella pneumoniae} held in the Advanced Research Laboratory. Distilled water produced none. GEN 5 remained the widest reference. Zones were measured inclusive of the disc. \taxa{E. coli} and \taxa{K. pneumoniae} gave larger means at C2 than at C1. \taxa{P. aeruginosa} gave its largest extract mean at C1. The aqueous extract is therefore active \invitro{} on these lawns, but activity is not a simple function of the two labelled strengths, and it is not a clinical result.

Future work should determine MIC and MBC in broth, preferably on Mueller--Hinton medium with identified strains. LC-MS or GC-MS of the aqueous concentrate would show which polar peaks are present. A wider concentration series would test whether the \taxa{P. aeruginosa} inversion repeats. Until those data exist, the present tables should be read as a preliminary aqueous screen of a Kashmir Botanical Garden collection, not as a treatment claim.

\FloatBarrier
